% BOOK-FILE-ID: f-ad0393fbde8e
% 解法の選択を、計算の前に確認する短い練習。
\subsection{操作を選ぶ練習}\label{節：操作を選ぶ練習}
一般項を計算する前に,置き換えた数列がどのような漸化式を満たすかを確かめる練習をしよう.
ここでは,引く・割る・\bookterm{logarithm}{対数}を取るという3つの操作を比べる.
どの操作を選べば既知の形になるのか,最後まで計算する前に見通せるようになることが目標である.

\noindent\bookblocklink{operation-choice}{problem}{answer}{【問題】}\par\nobreak
\begin{enumerate}[format=\bookitemlink{operation-choice}{problem}{answer}]
  \item $a_1=1$,\,$a_{n+1}=3a_n+2n+1$とする.
    定数を引くだけで付加項を消せるか.
    $g(n)=un+v$を引く場合の$u,v$を求め,\,$b_n=a_n-g(n)$が満たす漸化式と初項を示せ.
  \item $a_1=1$,\,$a_{n+1}=\frac{n+1}{n}a_n+(n+1)$とする.
    $b_n=\frac{a_n}{n}$と$c_n=\frac{a_n}{n(n+1)}$を候補とし,それぞれが満たす漸化式を求めよ.
    どちらで$a_n$の係数を$1$にできるか.
  \item $a_1=-2$,\,$a_{n+1}=2^{2n+1}a_n$とする.
    $b_n=\log_2a_n$は使えるか.
    \bookterm{logarithm}{対数}を使わず,\,$c_n=\frac{a_n}{2^{n^2}}$が満たす漸化式と初項を求めよ.
\end{enumerate}

\noindent\bookblocklink{operation-choice}{answer}{problem}{【解答と選択の理由】}\par\nobreak
\begin{enumerate}[format=\bookitemlink{operation-choice}{answer}{problem}]
  \item $a_n-\alpha$では付加項の$2n$が残る.
    $g(n+1)=3g(n)+2n+1$から
    \[
      u=3u+2,\qquad u+v=3v+1
    \]
    となり,\,$u=v=-1$である.
    従って$b_n=a_n+n+1$により$b_{n+1}=3b_n$,\,$b_1=3$を得る.
    引く量を定数から一次式へ変えたことで,付加項を消せた.

  \item それぞれの定義に元の漸化式を代入すると,
    \[
      b_{n+1}=b_n+1,\qquad b_1=1
    \]
    と
    \[
      c_{n+1}=\frac{n+1}{n+2}c_n+\frac{1}{n+2},\qquad c_1=\frac12
    \]
    を得る.\,$c_n$の式には係数$\frac{n+1}{n+2}$が残る.
    $h_n=n$なら
    \[
      \frac{h_{n+1}}{h_n}=\frac{n+1}{n}
    \]
    が元の係数と一致する.
    だからこの式で割る.
    $n\ge1$なので分母は0でなく,\,$b_n=n$から$a_n=n^2$も得られる.

  \item 初項が負で,毎回正の数を掛けるので,各項は負であり,\,$\log_2a_n$は実数として定義できない.
    一方,\,$2n+1=(n+1)^2-n^2$から
    \[
      c_{n+1}=c_n,\qquad c_1=-1
    \]
    となる.従って$a_n=-2^{n^2}$を得る.
    割る数$2^{n^2}$は常に0でないので,元の項が正である必要はない.
    初項が正の比較例は\bookproblemref{practice-basic}{42}{標準問42}へ.
\end{enumerate}

\paragraphbf{同じ式を2通りに解く}
第1問では,一次式$-n-1$を引くと等比型になり,\,$a_n=3^n-n-1$を得る.
一方,\,$c_n=\frac{a_n}{3^{n-1}}$と置いてもよい.この場合は
\[
  c_{n+1}-c_n=\frac{2n+1}{3^n}
\]
となり,残った和を計算して元へ戻す.
どちらも$n\ge1$で使える.ここでは一次式がすぐに決まり,引けば和の計算も要らない.
引く量が見つからなければ,係数から割る式を作れる.残る計算を比べて選ぶ.

\paragraphbf{両辺を割る方法の一般化}
第2問は\sectionref*{節：階比型}の操作を付加項がある場合へ広げたものだ.
一般に
\[
  a_{n+1}=p_na_n+q_n
\]
に対して,0にならない既知の\bookterm{sequence}{数列}$h_n$が$h_{n+1}=p_nh_n$を満たすなら,
\[
  b_n=\frac{a_n}{h_n},\qquad
  b_{n+1}-b_n=\frac{q_n}{h_{n+1}}
\]
となる.両辺を$h_{n+1}$で割った後は,\bookterm{difference}{階差}を書き下して足し合わせればよい.
$p_n\ne0$なら,\,$h_1=1$,\,$h_n=\prod_{k=1}^{n-1}p_k$を候補にできる(空積は$1$).
係数が0になる段階は別に調べる.残った和が簡単に計算できるとは限らない.
